% !TeX root = ../../main.tex
\section{案例五：CVE-2022-24877（Flux \texttt{kustomize\allowbreak-controller} 路径穿越）}

\subsection{漏洞详解与复现}

\noindent\hangindent=2em\hangafter=1 \textbf{漏洞背景}

Flux 是云原生 GitOps 交付体系中的常用组件，典型部署包含 \texttt{kustomize\allowbreak-controller} 等控制器，用于在集群中持续对齐 Git 仓库的期望状态。在该工作流中，控制器会拉取仓库内容并调用 Kustomize 解析 \texttt{kustomization\allowbreak.yaml}，进而生成并应用资源清单。由于 \texttt{kustomization\allowbreak.yaml} 以文件路径为核心配置对象，控制器必须将其中指定的资源、补丁与组件路径解析为实际文件系统访问。若路径规范化与边界检查不足，则恶意构造的 \texttt{kustomization\allowbreak.yaml} 可能触发路径穿越，使控制器读取超出预期工作目录的文件，从而造成敏感信息暴露，并在多租户部署中形成隔离边界破坏的风险。

\textbf{CVE 描述}：CVE-2022-24877 是 Flux \texttt{kustomize\allowbreak-controller} 的路径穿越漏洞。攻击者可通过恶意 \texttt{kustomization\allowbreak.yaml} 触发控制器读取其 Pod 文件系统中的非预期文件，从而暴露敏感数据；在多租户部署下，该能力可能进一步放大为权限提升风险。该漏洞已在 \texttt{kustomize\allowbreak-controller v0.24.0} 中修复，并随 \texttt{flux2 v0.29.0} 集成发布。工程侧常见缓解手段包括在 CI/CD 流水线中自动校验 \texttt{kustomization\allowbreak.yaml} 是否符合安全策略。

\textbf{主要成因分类}：输入校验缺陷 / 路径穿越（Path Traversal，路径规范化与目录边界检查不足）。

\textbf{攻击链条描述}：攻击链可概括为仓库内容被信任、路径解析越界以及敏感数据外泄。攻击者将恶意 \texttt{kustomization\allowbreak.yaml} 提交至可被控制器同步的仓库；控制器在构建阶段解析并访问其中指定的路径；若路径可越界访问控制器容器文件系统，则非预期文件内容可能被注入到构建产物、事件输出或后续流程中，造成敏感信息暴露，并在多租户环境中破坏隔离边界。

\noindent\hangindent=2em\hangafter=1 \textbf{漏洞复现（验证性复现）}

本案例复现采用\textbf{验证性复现}思路，重点验证控制器是否会在构建过程中读取超出预期目录边界的文件，并以构建日志与产物行为作为证据。复现的观测点包括：控制器在构建阶段出现与异常路径解析相关的错误或告警信息，或在严格受控的测试条件下出现构建产物包含非预期内容的迹象。为避免涉及敏感文件读取的可滥用细节，正文仅给出复现边界与证据判据。
完整的环境准备、受控验证步骤与证据判据见附录\ref{app:cve-2022-24877}。

\subsection{策略部署}

治理思路以前置校验与 GitOps 变更治理为核心。

\textbf{策略概述}：第一，在 CI/CD 或代码评审环节对 \texttt{kustomization\allowbreak.yaml} 进行策略化校验，阻断异常路径引用进入仓库。第二，在多租户环境中收敛对 Git 仓库写入与 Flux 资源变更的权限。第三，通过运行时硬化与最小权限配置降低异常路径被利用后的后果。

\textbf{策略配置与部署步骤}：CI 校验示例与权限收敛要点见附录\ref{app:cve-2022-24877}。

\subsection{三因素分析}

\noindent\textbf{有效性：}
CI 校验用于在提交阶段阻断异常路径引用进入仓库，从而显著降低触发概率。权限与变更治理用于限制恶意 \texttt{kustomization\allowbreak.yaml} 进入同步面，运行时硬化与最小权限配置则在多租户场景下缩小爆炸半径。

\noindent\textbf{无干扰性：}
CI 校验可能对历史仓库中不规范写法产生阻断，因此应采用先告警后强制的方式逐步推进，并提供例外与审批流程。多租户下的权限收敛可能影响自助交付效率，需要以标准化模板与可审计的变更机制替代高权限直写。运行时硬化通常对业务逻辑影响较小，但仍需验证其与现有控制器权限、文件访问行为及构建流程的兼容性。

\noindent\textbf{可部署性：}
上述措施以工程流程与平台能力为主，因而便于推广应用：CI 校验可作为流水线准入检查复用；多租户权限收敛与运行时硬化可通过 GitOps 模板化交付，实现持续改进。

就 CVE-2022-24877 而言，本文方法生成的策略通过 CI 流水线中的 \texttt{kustomization.yaml} 校验在提交阶段拦截异常路径引用，再配合多租户权限收敛与运行时硬化限制路径穿越被利用后的影响范围。该方案可作为补丁发布前的辅助缓解措施，最终修复仍依赖升级至 \texttt{kustomize-controller v0.24.0} 及以上版本。
